% Options for packages loaded elsewhere
\PassOptionsToPackage{unicode}{hyperref}
\PassOptionsToPackage{hyphens}{url}
\documentclass[
  11pt,
]{article}
\usepackage{xcolor}
\usepackage[left=3cm,right=3cm,top=2.5cm,bottom=2.5cm]{geometry}
\usepackage{amsmath,amssymb}
\setcounter{secnumdepth}{5}
\usepackage{iftex}
\ifPDFTeX
  \usepackage[T1]{fontenc}
  \usepackage[utf8]{inputenc}
  \usepackage{textcomp} % provide euro and other symbols
\else % if luatex or xetex
  \usepackage{unicode-math} % this also loads fontspec
  \defaultfontfeatures{Scale=MatchLowercase}
  \defaultfontfeatures[\rmfamily]{Ligatures=TeX,Scale=1}
\fi
\usepackage{lmodern}
\ifPDFTeX\else
  % xetex/luatex font selection
\fi
% Use upquote if available, for straight quotes in verbatim environments
\IfFileExists{upquote.sty}{\usepackage{upquote}}{}
\IfFileExists{microtype.sty}{% use microtype if available
  \usepackage[]{microtype}
  \UseMicrotypeSet[protrusion]{basicmath} % disable protrusion for tt fonts
}{}
\usepackage{setspace}
\makeatletter
\@ifundefined{KOMAClassName}{% if non-KOMA class
  \IfFileExists{parskip.sty}{%
    \usepackage{parskip}
  }{% else
    \setlength{\parindent}{0pt}
    \setlength{\parskip}{6pt plus 2pt minus 1pt}}
}{% if KOMA class
  \KOMAoptions{parskip=half}}
\makeatother
\usepackage{longtable,booktabs,array}
\usepackage{caption}
\captionsetup[table]{skip=6pt}
\newcounter{none} % for unnumbered tables
\usepackage{calc} % for calculating minipage widths
% Correct order of tables after \paragraph or \subparagraph
\usepackage{etoolbox}
\makeatletter
\patchcmd\longtable{\par}{\if@noskipsec\mbox{}\fi\par}{}{}
\makeatother
% Allow footnotes in longtable head/foot
\IfFileExists{footnotehyper.sty}{\usepackage{footnotehyper}}{\usepackage{footnote}}
\makesavenoteenv{longtable}
\usepackage{graphicx}
\makeatletter
\newsavebox\pandoc@box
\newcommand*\pandocbounded[1]{% scales image to fit in text height/width
  \sbox\pandoc@box{#1}%
  \Gscale@div\@tempa{\textheight}{\dimexpr\ht\pandoc@box+\dp\pandoc@box\relax}%
  \Gscale@div\@tempb{\linewidth}{\wd\pandoc@box}%
  \ifdim\@tempb\p@<\@tempa\p@\let\@tempa\@tempb\fi% select the smaller of both
  \ifdim\@tempa\p@<\p@\scalebox{\@tempa}{\usebox\pandoc@box}%
  \else\usebox{\pandoc@box}%
  \fi%
}
% Set default figure placement to htbp
\def\fps@figure{htbp}
\makeatother
% definitions for citeproc citations
\NewDocumentCommand\citeproctext{}{}
\NewDocumentCommand\citeproc{mm}{%
  \begingroup\def\citeproctext{#2}\cite{#1}\endgroup}
\makeatletter
 % allow citations to break across lines
 \let\@cite@ofmt\@firstofone
 % avoid brackets around text for \cite:
 \def\@biblabel#1{}
 \def\@cite#1#2{{#1\if@tempswa , #2\fi}}
\makeatother
\newlength{\cslhangindent}
\setlength{\cslhangindent}{1.5em}
\newlength{\csllabelwidth}
\setlength{\csllabelwidth}{3em}
\newenvironment{CSLReferences}[2] % #1 hanging-indent, #2 entry-spacing
 {\begin{list}{}{%
  \setlength{\itemindent}{0pt}
  \setlength{\leftmargin}{0pt}
  \setlength{\parsep}{0pt}
  % turn on hanging indent if param 1 is 1
  \ifodd #1
   \setlength{\leftmargin}{\cslhangindent}
   \setlength{\itemindent}{-1\cslhangindent}
  \fi
  % set entry spacing
  \setlength{\itemsep}{#2\baselineskip}}}
 {\end{list}}
\usepackage{calc}
\newcommand{\CSLBlock}[1]{\hfill\break\parbox[t]{\linewidth}{\strut\ignorespaces#1\strut}}
\newcommand{\CSLLeftMargin}[1]{\parbox[t]{\csllabelwidth}{\strut#1\strut}}
\newcommand{\CSLRightInline}[1]{\parbox[t]{\linewidth - \csllabelwidth}{\strut#1\strut}}
\newcommand{\CSLIndent}[1]{\hspace{\cslhangindent}#1}
\setlength{\emergencystretch}{3em} % prevent overfull lines
\providecommand{\tightlist}{%
  \setlength{\itemsep}{0pt}\setlength{\parskip}{0pt}}
% Paragraph scaffolding for manuscript revision.
%
% Two environments, presented in this order per paragraph:
%
%   rgt   -- author's rewrite (initially a placeholder)
%            small, italic, dark blue
%   orig  -- original Claude-authored draft, and until the rewrite
%            replaces it, the body text of the manuscript
%            normal size, gray
%
% The `bullets' environment (a boiled-down topic sentence plus bullet
% points, one per paragraph) is no longer used inside manuscripts:
% those blocks now live in a companion bullets.Rmd supplement in the
% same report directory, which typesets them as plain \textbf{} +
% itemize and needs no part of this file. The definition is retained
% only so that a stray block left in a document still compiles.
%
% This file is identical across every manuscript that includes it.
% It previously existed in two variants that swapped the sizes of
% rgt and orig, so the same markup rendered differently depending on
% which repository the manuscript lived in. The variant retained here
% keeps orig at full size because orig, not rgt, currently carries
% the prose in every manuscript in the pool.

\usepackage{xcolor}

\newcounter{pgraph}

\newenvironment{bullets}{%
  \stepcounter{pgraph}%
  \begin{quote}\small\sffamily\color[gray]{0.30}%
  {\normalfont\scriptsize\color[gray]{0.58}[\thepgraph]\enspace}}{%
  \end{quote}}

\newenvironment{rgt}{%
  \begin{quote}\small\itshape\color[RGB]{20,60,160}}{%
  \end{quote}}

\newenvironment{orig}{%
  \begin{quote}\normalsize\color[gray]{0.55}}{%
  \end{quote}}

% backward compat alias
\newenvironment{claudecode}{\begin{orig}}{\end{orig}}
%% tools/stamp.tex
%% Provenance footer for PDFs rendered from this project.
%% Vendored by zzcollab; this file is not project-specific. The
%% footer prints the source document and the time of rendering.
%% The git version is defined here too, and so is available to a
%% project that wants it on the page, but it is left off the
%% default footer: it already names the staged copy in share/ and
%% appears in its manifest row. All values are supplied at render
%% time by tools/stamp-render.R, which writes a small generated
%% file that redefines the macros below. The \providecommand
%% defaults keep a bare LaTeX compile working when the document is
%% built without the wrapper.
\usepackage{fancyhdr}
\usepackage{xcolor}
\providecommand{\stampsource}{\detokenize{(source not recorded)}}
\providecommand{\stamptime}{(time not recorded)}
\providecommand{\stampversion}{\detokenize{(version not recorded)}}
\setlength{\footskip}{40pt}
\newcommand{\stampfootcontent}{%
  \footnotesize\color{gray}%
  \parbox[b]{\textwidth}{%
    \stamptime\hfill\thepage\\[2pt]%
    \ttfamily\raggedright\stampsource}}
\pagestyle{fancy}
\fancyhf{}
\fancyfoot[C]{\stampfootcontent}
\renewcommand{\headrulewidth}{0pt}
\renewcommand{\footrulewidth}{0.4pt}
%% The title page uses the `plain` page style; redefine it so the
%% stamp appears there too.
\fancypagestyle{plain}{%
  \fancyhf{}%
  \fancyfoot[C]{\stampfootcontent}%
  \renewcommand{\headrulewidth}{0pt}%
  \renewcommand{\footrulewidth}{0.4pt}}
\renewcommand{\stampsource}{\textasciitilde{}/\allowbreak{}prj/\allowbreak{}res/\allowbreak{}36-pmsimstats-ng/\allowbreak{}pmsimstats-ng/\allowbreak{}analysis/\allowbreak{}report/\allowbreak{}12-nof1-design-efficiency/\allowbreak{}report.Rmd}
\renewcommand{\stamptime}{2026-08-20 18:07 PDT}
\renewcommand{\stampversion}{\detokenize{e65895c}}
\usepackage{amsmath}
\usepackage{amssymb}
\usepackage{booktabs}
\usepackage{longtable}
\usepackage{float}
\usepackage[mathlines]{lineno}
\linenumbers
\usepackage{bookmark}
\IfFileExists{xurl.sty}{\usepackage{xurl}}{} % add URL line breaks if available
\urlstyle{same}
% fallback for those not using the hyperref driver hyperxmp:
\makeatletter
\@ifundefined{xmpquote}{\newcommand{\xmpquote}[1]{#1}}{}
\makeatother
\hypersetup{
  pdftitle={A closed-form approximation to the biomarker-treatment interaction estimator in aggregated N-of-1 trials{,} and what it implies for design efficiency under a mandatory open-label lead-in},
  pdfauthor={pmsimstats team},
  hidelinks,
  pdfcreator={LaTeX via pandoc}}

\title{A closed-form approximation to the biomarker-treatment interaction estimator in aggregated N-of-1 trials, and what it implies for design efficiency under a mandatory open-label lead-in}
\author{pmsimstats team}
\date{August 20, 2026}

\begin{document}
\maketitle

{
\setcounter{tocdepth}{2}
\tableofcontents
}
\setstretch{1.1}
\textbf{Short communication.}

\section{Introduction}\label{introduction}

\begin{rgt}
Lorem ipsum dolor sit amet, consectetur adipiscing elit, sed do eiusmod tempor incididunt ut labore et dolore magna aliqua.

\end{rgt}

\begin{orig}
Every design question this project's companion manuscripts have
addressed --- how many randomization paths, what sample size, how
long to wait before a discontinuation-phase assessment, whether a
brief crossover tail is worth its extra visits --- has so far been
answered by Monte Carlo simulation\textsuperscript{\citeproc{ref-pmsimstats-paper01}{1}--\citeproc{ref-pmsimstats-paper11}{5}}. That is appropriate for
the compendium's headline estimand, which combines a Gompertz
response trajectory, a dual-channel biomarker-moderation
architecture, exponential or Weibull carryover, and a continuous-time
AR(1) mixed model into a system with no obvious closed form. But it
leaves a practical gap: a trialist choosing among two or three
concrete design options today cannot get even an order-of-magnitude
power estimate without running a fresh simulation, in contrast to the
treatment main effect, for which Senn's and Zucker's classical
variance formulas\textsuperscript{\citeproc{ref-sennSampleSizeConsiderations2019}{6},\citeproc{ref-zucker2010}{7}}
have long supplied exactly this kind of design intuition, and in
contrast to the dichotomized-biomarker case, for which the companion
test-procedure manuscript derives a closed-form repeated-measures
ANOVA power formula\textsuperscript{\citeproc{ref-pmsimstats-paper10}{8}}. Neither covers the
compendium's actual estimand: a continuous biomarker, a continuous
exposure predictor \(D_{bc}\), and a mixed model fit by
\texttt{nlme::lme} with \texttt{corCAR1} residual correlation.

\end{orig}

\begin{rgt}
Lorem ipsum dolor sit amet, consectetur adipiscing elit, sed do eiusmod tempor incididunt ut labore et dolore magna aliqua.

\end{rgt}

\begin{orig}
This short communication was motivated by a concrete instance of the
gap: a trial whose first phase must be open-label, for recruitment
reasons that are not themselves a design choice. Given that fixed
constraint, what should the analyst prioritize among the remaining
design parameters to make the biomarker-treatment interaction test as
efficient as possible? We first give a qualitative answer, drawn
directly from the compendium's existing simulation results. We then
ask whether a genuine closed-form approximation to the interaction
estimator's sampling variance and power can be derived, and, if so,
how far it can be trusted. We derive such an approximation from two
components: a within-subject two-group contrast reduction of the
analysis model, and the exact variance of a mean of AR(1)-correlated,
unequally spaced observations. We validate both components against
Monte Carlo output already on disk, produced for a companion
manuscript, rather than commissioning a new simulation study --- the
appropriate scope for a short communication. We report the validation
honestly: the variance and power approximation performs well: the
attenuation-factor component captures the correct qualitative ranking
of designs by carryover vulnerability but not its magnitude. We close
with quantitative, closed-form-grounded design guidance for the
open-label-first scenario.

\end{orig}

\section{What the compendium already establishes}\label{what-the-compendium-already-establishes}

\begin{rgt}
Lorem ipsum dolor sit amet, consectetur adipiscing elit, sed do eiusmod tempor incididunt ut labore et dolore magna aliqua.

\end{rgt}

\begin{orig}
Before deriving anything new, it is worth stating precisely what the
existing simulation evidence already supports, since a closed form is
only useful to the extent that it reproduces conclusions already
established with more care. Four results from the compendium bear
directly on the open-label-first design question.

First, design topology after the open-label phase is the largest
lever available. At matched total sample size \(N=70\), the Hybrid
design (open-label, then blinded discontinuation, then a brief
crossover; four randomization paths) dominates open-label-plus-
blinded-discontinuation alone (OL+BDC; two paths) in absolute power
at every carryover half-life tested, \(0.842\) versus \(0.751\) at
\(t_{1/2}=0\) and \(0.784\) versus \(0.730\) at \(t_{1/2}=1.0\) week\textsuperscript{\citeproc{ref-pmsimstats-paper01}{1}}. Adding the crossover tail is worth its extra
visits.

Second, matching total sample size \(N\) across designs being compared
is not optional bookkeeping: it changed the compendium's own
conclusions twice. An unmatched comparison inflated Hybrid's apparent
power ceiling in the architecture manuscript and manufactured an
apparent dropout-robustness advantage for Hybrid in the
informative-dropout manuscript that a matched re-run showed was a
ceiling effect\textsuperscript{\citeproc{ref-pmsimstats-paper01}{1},\citeproc{ref-pmsimstats-paper09}{4}}.

Third, the carryover decay's functional \emph{shape} (exponential versus
Weibull, across shape parameters \(k \in \{0.25, 0.5, 2, 4\}\)) barely
moves the compendium's substantive conclusions; the \emph{analysis
specification} --- whether the exposure predictor is a binary
indicator, an exposure-weighted continuous decay, or an
AIC-selected-half-life variant --- matters considerably more\textsuperscript{\citeproc{ref-pmsimstats-paper02}{2}}.

Fourth, misspecifying the residual correlation structure (fitting an
exchangeable or naive model when the truth is AR(1)) is a larger
anti-conservatism risk than misspecifying carryover decay shape\textsuperscript{\citeproc{ref-pmsimstats-paper10}{8}}.

These four results already answer ``what to prioritize'' at a
qualitative level. What they do not supply is a way to move from a
verbal ranking to a number: how much does adding the crossover tail
buy, quantitatively, as a function of the design's own visit
schedule, without running a new simulation for every candidate
design? That is the question the rest of this communication addresses.

\end{orig}

\section{A closed-form approximation}\label{a-closed-form-approximation}

\subsection{Setup}\label{setup}

\begin{rgt}
Lorem ipsum dolor sit amet, consectetur adipiscing elit, sed do eiusmod tempor incididunt ut labore et dolore magna aliqua.

\end{rgt}

\begin{orig}
We work under the mean-moderation architecture\textsuperscript{\citeproc{ref-pmsimstats-paper01}{1}}, in which the biomarker \(B_i\) scales the
treatment effect directly in the outcome's conditional mean rather
than through a covariance channel, and we approximate the analysis
specification E1 (binary on-drug indicator \(D_{it}\), carryover
unmodeled at the analysis stage;\textsuperscript{\citeproc{ref-pmsimstats-paper02}{2}}) by a simple
two-group within-subject contrast. For participant \(i\) with \(k_{on}\)
on-drug and \(k_{off}\) off-drug assessment timepoints at fixed
calendar times, define
\begin{equation}
\Delta_i = \overline{Y}_{i}^{\,on} - \overline{Y}_{i}^{\,off},
\label{eq:contrast}
\end{equation}
the mean on-drug outcome minus the mean off-drug outcome. Because the
participant random intercept \(u_i\) is common to both means, it
cancels in \(\Delta_i\) exactly (no random slope is present in the
compendium's baseline model), leaving
\begin{equation}
\Delta_i = \beta_1\big(1 + \beta_{bm} B_i\big)\cdot\bar D_{bc,i}^{\,on}
  - \beta_1\big(1+\beta_{bm}B_i\big)\cdot\bar D_{bc,i}^{\,off} + \eta_i,
\label{eq:delta-decomp}
\end{equation}
where \(\bar D_{bc,i}^{\,on}=1\) by construction and \(\bar
D_{bc,i}^{\,off} = \overline{D_{bc}}_i^{\,off}\) is the mean of the
continuous exposure-decayed indicator across the off-drug assessment
times, which is less than one whenever residual carryover has not
fully decayed by the time of assessment, and \(\eta_i\) is the
difference of the two residual means. A second-stage ordinary least
squares regression of \(\Delta_i\) on the standardized biomarker \(b_i\)
then recovers the interaction coefficient as its slope,
\begin{equation}
\Delta_i = \gamma_0 + \beta_{bm:D}\, b_i + \eta_i, \qquad
  \beta_{bm:D} = \beta_1\,\beta_{bm}\cdot A_i,
\label{eq:e9-like}
\end{equation}
where \(A_i = 1 - \overline{D_{bc}}_i^{\,off}\) is a per-participant
\textbf{carryover attenuation factor}. This construction is exactly the
two-stage paired-difference specification E9 of the companion
carryover manuscript\textsuperscript{\citeproc{ref-pmsimstats-paper02}{2}}, specialized to a strict
on/off split rather than E9's general on/off means, and is offered
here as an approximation to E1 (which is fit by the full mixed model
on all timepoints jointly, not by this two-group reduction) rather
than as an exact restatement of it.

\end{orig}

\subsection{Sampling variance under AR(1) correlation}\label{sampling-variance-under-ar1-correlation}

\begin{rgt}
Lorem ipsum dolor sit amet, consectetur adipiscing elit, sed do eiusmod tempor incididunt ut labore et dolore magna aliqua.

\end{rgt}

\begin{orig}
The variance of the mean of \(k\) observations drawn from a stationary
process with autocovariance \(\gamma(h) = \sigma^2\rho^{|h|}\) at
arbitrary, possibly unequally spaced times \(t_1,\dots,t_k\) has the
exact closed form
\begin{equation}
V(\mathbf{t}, \sigma^2, \rho) = \frac{\sigma^2}{k^2}
  \sum_{j=1}^{k}\sum_{l=1}^{k} \rho^{|t_j - t_l|}.
\label{eq:var-mean-ar1}
\end{equation}
Applying Equation \eqref{eq:var-mean-ar1} separately to a
participant's on-drug and off-drug visit times and summing (treating
the two period means as approximately independent, an approximation
we return to in the Discussion) gives the variance of the
within-subject contrast,
\begin{equation}
\mathrm{Var}(\eta_i) \approx V(\mathbf{t}^{on}, \sigma^2, \rho) +
  V(\mathbf{t}^{off}, \sigma^2, \rho) \;\equiv\; v(\mathrm{design},
  \rho),
\label{eq:v-sum}
\end{equation}
a quantity computable directly from a design's visit schedule and an
assumed AR(1) parameter \(\rho\), with no simulation required. Pooling
across \(N\) participants (assuming a standardized biomarker, so
\(\mathrm{Var}(b_i)=1\)) and, for a multi-path design, averaging
\(v(\mathrm{design},\rho)\) across the design's randomization paths,
the closed-form approximation to the interaction estimator's sampling
variance is
\begin{equation}
\widehat{\mathrm{Var}}\big(\hat\beta_{bm:D}\big) \approx
  \frac{\sigma^2\, v_1(\mathrm{design},\rho)}{N},
\label{eq:var-closed-form}
\end{equation}
where \(v_1\) is \(v(\mathrm{design},\rho)\) evaluated at \(\sigma^2=1\), so
that \(\sigma^2\) enters as a single free scale parameter separating
design-and-correlation structure (\(v_1\), closed-form) from residual
noise magnitude (\(\sigma^2\), an empirical quantity like any other
nuisance variance component).

\end{orig}

\subsection{Predicted power}\label{predicted-power}

\begin{rgt}
Lorem ipsum dolor sit amet, consectetur adipiscing elit, sed do eiusmod tempor incididunt ut labore et dolore magna aliqua.

\end{rgt}

\begin{orig}
Given Equations \eqref{eq:e9-like} and \eqref{eq:var-closed-form}, a
Wald-normal approximation to power at level \(\alpha\) follows
immediately, in the same spirit as the general four-channel
decomposition of the companion calibration manuscript\textsuperscript{\citeproc{ref-pmsimstats-paper10}{8}, its \(\kappa=1,\nu=\infty\) special case}:
\begin{equation}
z = \frac{|\beta_{bm:D}\cdot A(t_{1/2})|}
  {\sqrt{\sigma^2\, v_1(\mathrm{design},\rho)/N}}, \qquad
  \hat\pi \approx \Phi(z - z_{1-\alpha/2}) + \Phi(-z - z_{1-\alpha/2}),
\label{eq:power-closed-form}
\end{equation}
with \(A(t_{1/2})\) the design-averaged attenuation factor of
Equation \eqref{eq:e9-like} and \(\beta_{bm:D}\) the interaction
coefficient's magnitude in the absence of carryover. Every quantity
on the right-hand side except \(\sigma^2\), \(\rho\), and \(\beta_{bm:D}\)
itself is computable directly from the candidate design's visit
schedule and a candidate carryover half-life, with no simulation.

\end{orig}

\section{Validation against existing Monte Carlo output}\label{validation-against-existing-monte-carlo-output}

\subsection{Data and procedure}\label{data-and-procedure}

\begin{rgt}
Lorem ipsum dolor sit amet, consectetur adipiscing elit, sed do eiusmod tempor incididunt ut labore et dolore magna aliqua.

\end{rgt}

\begin{orig}
We validate Equations \eqref{eq:var-closed-form} and
\eqref{eq:power-closed-form} against Monte Carlo output already
produced for the dual-channel-architecture companion manuscript\textsuperscript{\citeproc{ref-pmsimstats-paper11}{5}}, rather than running a new simulation. That
manuscript's driver
(\texttt{analysis/scripts/dgp-combined/01-run-combined-factorial.R})
generates data under a combined architecture with independent
mean-channel and covariance-channel weights \(c_{bm,A}\) and
\(c_{bm,B}\); its \((c_{bm,A}=0.45,\, c_{bm,B}=0)\) boundary is, by
construction, exactly the pure mean-moderation architecture this
communication's derivation assumes. Its stored summary
(\texttt{analysis/data/dgp-combined/01-combined-summary.rds}) already
includes matched total sample sizes \(N \in \{35, 70\}\), three designs
(crossover, CO; open-label-plus-blinded-discontinuation, OL+BDC;
Hybrid), three carryover half-lives \(t_{1/2}\in\{0, 0.5, 1.0\}\)
weeks, and analysis specification E1, at \(1{,}000\) replicates per
cell. We compute \(v_1(\mathrm{design}, \rho{=}0.7)\) and
\(A(\mathrm{design}, t_{1/2})\) directly from the same visit schedules
the driver itself uses\textsuperscript{\citeproc{ref-pmsimstats-paper02}{2(sec2)} for the design
definitions, reproduced with driver-matching timepoints in
\texttt{analysis/scripts/nof1-design-efficiency/01-closed-form-validation.R}},
fit the single free scale parameter \(\sigma^2\) by ordinary least
squares of \(\widehat{\mathrm{sd}}(\hat\beta)^2 \cdot N\) on \(v_1\)
across all 18 design-by-half-life-by-\(N\) cells, and then compare the
resulting closed-form power predictions to the cells' simulated
power. \(\rho=0.7\) is the AR(1) parameter the driver itself uses by
default and is not separately estimated.

\end{orig}

\subsection{The variance-scaling component performs well}\label{the-variance-scaling-component-performs-well}

\begin{rgt}
Lorem ipsum dolor sit amet, consectetur adipiscing elit, sed do eiusmod tempor incididunt ut labore et dolore magna aliqua.

\end{rgt}

\begin{orig}
With a single fitted scale parameter (\(\hat\sigma^2 = 0.667\)), the
closed form explains \(R^2 = 0.948\) of the variation in
\(\widehat{\mathrm{sd}}(\hat\beta)^2 \cdot N\) across the 18 cells
(Table \ref{tab:variance-fit}). Two structural predictions of
Equation \eqref{eq:var-closed-form} are confirmed directly: the
sampling variance scales as \(1/N\) almost exactly (the CO design's
empirical standard deviation ratio between \(N=35\) and \(N=70\) is
\(1.439\), against a predicted \(\sqrt{70/35}=1.414\)), and the design
ranking by \(v_1\) (OL+BDC highest, Hybrid intermediate, CO lowest, at
fixed \(\rho\)) is reproduced in direction, though not exactly in
magnitude, by the empirical standard deviations. A systematic
residual pattern is visible and is reported rather than concealed:
the closed form under-predicts CO's empirical standard deviation by
roughly 15-20\% and over-predicts Hybrid's and OL+BDC's by roughly
5-10\%. The most likely source is Equation \eqref{eq:v-sum}`s
independence approximation between a participant's on-drug and
off-drug period means, which is exact only when no AR(1) correlation
crosses the on/off boundary; CO's two long, contiguous four-visit
blocks per path plausibly violate that approximation more than the
other designs' more fragmented schedules.

\end{orig}

\begin{table}[H]
\centering
\caption{Closed-form sampling-standard-deviation predictions against
Monte Carlo output, at $\hat\sigma^2 = 0.667$ (fitted once, across
all 18 cells). $\mathrm{sd}(\hat\beta)$ is the empirical standard
deviation of the fitted \texttt{bm:Dbc} coefficient across 1{,}000
replicates; $\widehat{\mathrm{sd}}(\hat\beta)$ is the closed-form
prediction of Equation \ref{eq:var-closed-form}.}
\label{tab:variance-fit}
\begin{tabular}{lrrrrr}
\toprule
Design & $N$ & $t_{1/2}$ (wk) & $\mathrm{sd}(\hat\beta)$ &
  $\widehat{\mathrm{sd}}(\hat\beta)$ & Ratio \\
\midrule
CO     & 35 & 0.0 & 0.159 & 0.132 & 1.21 \\
CO     & 70 & 0.0 & 0.110 & 0.093 & 1.18 \\
Hybrid & 35 & 0.0 & 0.123 & 0.134 & 0.92 \\
Hybrid & 70 & 0.0 & 0.084 & 0.095 & 0.89 \\
OL+BDC & 35 & 0.0 & 0.138 & 0.147 & 0.93 \\
OL+BDC & 70 & 0.0 & 0.096 & 0.104 & 0.92 \\
CO     & 35 & 1.0 & 0.153 & 0.132 & 1.16 \\
CO     & 70 & 1.0 & 0.108 & 0.093 & 1.16 \\
Hybrid & 35 & 1.0 & 0.120 & 0.134 & 0.89 \\
Hybrid & 70 & 1.0 & 0.087 & 0.095 & 0.92 \\
OL+BDC & 35 & 1.0 & 0.139 & 0.147 & 0.94 \\
OL+BDC & 70 & 1.0 & 0.096 & 0.104 & 0.92 \\
\bottomrule
\end{tabular}
\end{table}

\begin{rgt}
Lorem ipsum dolor sit amet, consectetur adipiscing elit, sed do eiusmod tempor incididunt ut labore et dolore magna aliqua.

\end{rgt}

\begin{orig}
The \(t_{1/2}=0.5\) rows, omitted from Table \ref{tab:variance-fit}
for space, follow the same pattern and are included in the full
18-row output of
\texttt{analysis/data/nof1-design-efficiency/closed-form-validation.csv}.

\end{orig}

\subsection{Power prediction}\label{power-prediction}

\begin{rgt}
Lorem ipsum dolor sit amet, consectetur adipiscing elit, sed do eiusmod tempor incididunt ut labore et dolore magna aliqua.

\end{rgt}

\begin{orig}
Using the fitted \(\hat\sigma^2\) and, for each cell, that cell's own
empirical \(\overline{\hat\beta}\) as the effective true effect in
Equation \eqref{eq:power-closed-form} --- a choice discussed in
\S\ref{limitations} --- the closed form predicts simulated power with
a mean absolute error of \(0.049\), a root-mean-square error of
\(0.058\), and a correlation of \(0.935\) across the 18 cells (Table
\ref{tab:power-fit} gives the \(t_{1/2}=0\) and \(t_{1/2}=1.0\) rows in
full). For a first-pass closed form validated against a single
architecture at one \(\rho\) value, this is a substantively useful
level of agreement: a trialist screening candidate designs would
reach the same qualitative decision from the closed form as from the
full Monte Carlo simulation in the great majority of cells inspected
here.

\end{orig}

\begin{table}[H]
\centering
\caption{Closed-form power predictions against simulated power
(1{,}000 replicates per cell, $\alpha=0.05$ two-sided).}
\label{tab:power-fit}
\begin{tabular}{lrrrrr}
\toprule
Design & $N$ & $t_{1/2}$ (wk) & Power (simulated) &
  Power (closed form) & Error \\
\midrule
CO     & 35 & 0.0 & 0.352 & 0.421 & $+0.069$ \\
CO     & 70 & 0.0 & 0.595 & 0.690 & $+0.095$ \\
Hybrid & 35 & 0.0 & 0.452 & 0.426 & $-0.026$ \\
Hybrid & 70 & 0.0 & 0.738 & 0.703 & $-0.035$ \\
OL+BDC & 35 & 0.0 & 0.335 & 0.367 & $+0.032$ \\
OL+BDC & 70 & 0.0 & 0.612 & 0.618 & $+0.006$ \\
CO     & 35 & 1.0 & 0.323 & 0.409 & $+0.086$ \\
CO     & 70 & 1.0 & 0.611 & 0.716 & $+0.105$ \\
Hybrid & 35 & 1.0 & 0.456 & 0.418 & $-0.038$ \\
Hybrid & 70 & 1.0 & 0.729 & 0.688 & $-0.041$ \\
OL+BDC & 35 & 1.0 & 0.327 & 0.357 & $+0.030$ \\
OL+BDC & 70 & 1.0 & 0.577 & 0.596 & $+0.019$ \\
\bottomrule
\end{tabular}
\end{table}

\begin{rgt}
Lorem ipsum dolor sit amet, consectetur adipiscing elit, sed do eiusmod tempor incididunt ut labore et dolore magna aliqua.

\end{rgt}

\begin{orig}
The closed form correctly identifies Hybrid as the best-powered
design at both sample sizes and both half-lives shown, matching the
compendium's own headline finding\textsuperscript{\citeproc{ref-pmsimstats-paper01}{1}}. It does not
correctly order CO against OL+BDC: the closed form ranks CO above
OL+BDC at \(N=70\) (predicted power \(0.690\) versus \(0.618\) at
\(t_{1/2}=0\)), while the simulation ranks them the other way (\(0.595\)
versus \(0.612\)). This is a genuine limitation, not a rounding
artefact, and is reported as such rather than smoothed over: the
closed form is reliable for identifying the best design among
candidates that differ substantially in \(v_1\), and unreliable for
close calls between designs whose \(v_1\) values are similar.

\end{orig}

\subsection{The attenuation factor: correct ranking, overstated magnitude}\label{the-attenuation-factor-correct-ranking-overstated-magnitude}

\begin{rgt}
Lorem ipsum dolor sit amet, consectetur adipiscing elit, sed do eiusmod tempor incididunt ut labore et dolore magna aliqua.

\end{rgt}

\begin{orig}
Table \ref{tab:attenuation} compares the closed-form attenuation
factor \(A(t_{1/2})\), computed purely from each design's visit
schedule, to the actual relative change in the fitted mixed-model
coefficient's magnitude, \(|\overline{\hat\beta}(t_{1/2}{=}1.0)| \,/\,
|\overline{\hat\beta}(t_{1/2}{=}0)|\), from the same Monte Carlo cells.
The direction is right in every case, and the ranking across designs
is right: CO's off-drug visits are, on one randomization path, a
pre-drug baseline period that carryover cannot contaminate at all, so
CO shows the least closed-form attenuation of the three designs,
exactly matching the compendium's own qualitative finding that CO's
carryover-induced power loss is the smallest of the designs compared\textsuperscript{\citeproc{ref-pmsimstats-paper01}{1}}. But the \emph{magnitude} is badly overstated
throughout: the closed form predicts a \(24.5\)-percentage-point
attenuation for Hybrid between \(t_{1/2}=0\) and \(t_{1/2}=1.0\) weeks,
while the fitted coefficient's magnitude falls by only about \(1.7\%\)
over the same range (\(N=70\): \(|{-}0.2358|\) to \(|{-}0.2317|\)).

\end{orig}

\begin{table}[H]
\centering
\caption{Closed-form attenuation factor versus the actual relative
change in the fitted interaction coefficient's magnitude, $N=70$.}
\label{tab:attenuation}
\begin{tabular}{lrr}
\toprule
Design & $A(t_{1/2}{=}1.0)$ (closed form) &
  $|\bar{\hat\beta}(1.0)/\bar{\hat\beta}(0)|$ (simulated) \\
\midrule
CO     & 0.973 & 1.031 \\
Hybrid & 0.755 & 0.983 \\
OL+BDC & 0.667 & 0.975 \\
\bottomrule
\end{tabular}
\end{table}

\begin{rgt}
Lorem ipsum dolor sit amet, consectetur adipiscing elit, sed do eiusmod tempor incididunt ut labore et dolore magna aliqua.

\end{rgt}

\begin{orig}
The explanation is structural, not a fitting failure. The closed
form's attenuation factor comes from a crude two-group mean split
that discards each observation's individual timing; the fitted mixed
model instead includes an explicit time covariate \(t\) alongside
\texttt{bm:Dbc}, and recovers gradient information within the off-drug
period (how the outcome trends as \(D_{bc}\) decays across successive
off-drug visits) that the two-group contrast simply averages away.
This is a genuine, useful negative result: it means the attenuation
factor of Equation \eqref{eq:e9-like} should be read as a \emph{design
screening heuristic} --- it will correctly tell an analyst which of
two designs loses more information to a given carryover half-life
--- and not as a \emph{quantitative bias estimate} for the coefficient a
properly specified mixed-model analysis will actually report.

\end{orig}

\section{Design guidance for a mandatory open-label lead-in}\label{design-guidance-for-a-mandatory-open-label-lead-in}

\begin{rgt}
Lorem ipsum dolor sit amet, consectetur adipiscing elit, sed do eiusmod tempor incididunt ut labore et dolore magna aliqua.

\end{rgt}

\begin{orig}
Given a fixed open-label first phase, the validated components above
translate the compendium's qualitative guidance into a design
procedure. First, compute \(v_1(\mathrm{design}, \rho)\) from
Equation \eqref{eq:v-sum} for every design under consideration, using
each candidate's actual visit schedule; this can be done for a
design that has never been simulated, in seconds, without touching
the simulation pipeline. Second, use \(v_1\)'s ranking to eliminate
designs that are clearly dominated (a large \(v_1\) relative to the
alternatives) and to identify a small shortlist worth a full Monte
Carlo comparison; Table \ref{tab:variance-fit}'s validation shows
this ranking is trustworthy for well-separated candidates and should
not be trusted to resolve close calls, which still warrant
simulation. Third, use \(A(t_{1/2})\) from Equation \eqref{eq:e9-like}
only as a relative screen across designs at a shared assumed
half-life, not as a quantitative power adjustment; §4.4 shows it
substantially overstates the actual attenuation a \texttt{corCAR1} mixed
model with a time covariate will recover from. Fourth, since adding
visits after the mandatory open-label phase is the largest lever in
\(v_1\)'s design-dependence (Table \ref{tab:variance-fit}), and since
the compendium's own matched-\(N\) comparison shows a brief crossover
tail after blinded discontinuation is worth its cost in the one
design pairing checked here\textsuperscript{\citeproc{ref-pmsimstats-paper01}{1}}, a Hybrid-style
continuation of the mandatory open-label phase should be the default
starting point for the design shortlist, to be displaced only by a
full simulation showing a specific alternative outperforms it for
the trial's actual parameters.

\end{orig}

\section{Limitations}\label{limitations}

\begin{rgt}
Lorem ipsum dolor sit amet, consectetur adipiscing elit, sed do eiusmod tempor incididunt ut labore et dolore magna aliqua.

\end{rgt}

\begin{orig}
Several scope limitations bound how far this closed form should be
trusted. It is derived and validated under the mean-moderation
architecture only; the covariance-moderation architecture's
interaction lives in \(\partial\,\mathrm{Cor}(BM,BR)/\partial D_{bc}\)
rather than in the conditional mean\textsuperscript{\citeproc{ref-pmsimstats-paper01}{1}}, and the
two-group contrast reduction of Equation \eqref{eq:contrast} does not
apply to it without a separate derivation. The validation uses the
Monte Carlo cells' own empirical \(\overline{\hat\beta}\) as the
effective true effect in Equation \eqref{eq:power-closed-form} rather
than the architecture's nominal calibration target
\(c_{bm,A}\cdot\sigma_{BR}\): the two differ by roughly an order of
magnitude in the data inspected here (\(3.6\) nominal against
approximately \(0.23\) observed), a discrepancy this communication
does not resolve and flags as an open calibration question for the
companion architecture manuscripts rather than papering over with an
ad hoc rescaling. The variance formula's on/off independence
approximation (Equation \eqref{eq:v-sum}) is shown, by its residual
pattern in Table \ref{tab:variance-fit}, to be imperfect for designs
with long contiguous within-phase blocks. Validation is against a
single \(\rho\) value (\(0.7\), the driver default) and a single pair of
sample sizes; whether the closed form's accuracy holds at other
correlation strengths or substantially larger or smaller \(N\) is not
established here. Finally, the closed form omits the placebo-belief
and natural-history components of the outcome decomposition
entirely\textsuperscript{\citeproc{ref-pmsimstats-paper06}{9}}; a biomarker correlated with placebo
responsiveness will bias the interaction estimate in a way no version
of this closed form, restricted to the pharmacological channel alone,
can detect.

\end{orig}

\section{Conclusion}\label{conclusion}

\begin{rgt}
Lorem ipsum dolor sit amet, consectetur adipiscing elit, sed do eiusmod tempor incididunt ut labore et dolore magna aliqua.

\end{rgt}

\begin{orig}
A closed-form approximation to the biomarker-treatment interaction
estimator's sampling variance, built from a within-subject two-group
contrast and the exact variance of a mean of AR(1)-correlated
observations, reproduces 95\% of the variation in the interaction
estimator's Monte Carlo sampling variance and predicts simulated
power to within a mean absolute error of 0.05, using a single fitted
nuisance-variance parameter and no new simulation. A companion
attenuation factor correctly ranks trial designs by their
vulnerability to carryover but should not be read as a quantitative
bias estimate, because it discards timing information a properly
specified mixed-model analysis recovers. Together these give a
concrete, quantitative starting point for the design question that
motivated this communication --- what to prioritize given a mandatory
open-label first phase --- without displacing the full Monte Carlo
comparison the compendium's other manuscripts continue to supply for
final design decisions.

\end{orig}

\section{Reproducibility}\label{reproducibility}

\begin{orig}
No new Monte Carlo simulation was run for this manuscript. All
numbers reported here are produced by
\texttt{analysis/scripts/nof1-design-efficiency/01-closed-form-validation.R},
which reads the existing output of
\texttt{analysis/scripts/dgp-combined/01-run-combined-factorial.R}
(\texttt{analysis/data/dgp-combined/01-combined-summary.rds}) and writes
\texttt{analysis/data/nof1-design-efficiency/closed-form-validation.\{rds,csv\}}.

\end{orig}

\section*{References}\label{references}
\addcontentsline{toc}{section}{References}

\protect\phantomsection\label{refs}
\begin{CSLReferences}{0}{1}
\bibitem[\citeproctext]{ref-pmsimstats-paper01}
\CSLLeftMargin{1. }%
\CSLRightInline{{pmsimstats team}. Two architectures for simulating biomarker-treatment interactions: Implications for statistical power under carryover. Published online 2026.}

\bibitem[\citeproctext]{ref-pmsimstats-paper02}
\CSLLeftMargin{2. }%
\CSLRightInline{{pmsimstats team}. Carryover representation and inference for biomarker-treatment interaction tests in aggregated {N}-of-1 trials under covariance moderation: A factorial simulation comparison of nine analysis specifications. Published online 2026.}

\bibitem[\citeproctext]{ref-pmsimstats-paper05}
\CSLLeftMargin{3. }%
\CSLRightInline{{pmsimstats team}. Design sensitivity of aggregated {N}-of-1 trials: A twelve-sweep simulation study of the treatment main effect. Published online 2026.}

\bibitem[\citeproctext]{ref-pmsimstats-paper09}
\CSLLeftMargin{4. }%
\CSLRightInline{{pmsimstats team}. Informative dropout by trial design in aggregated {N}-of-1 biomarker-treatment interaction trials. Published online 2026.}

\bibitem[\citeproctext]{ref-pmsimstats-paper11}
\CSLLeftMargin{5. }%
\CSLRightInline{{pmsimstats team}. A combined data-generating architecture for biomarker-treatment interactions: Channel super-additivity and mixing-weight sensitivity in aggregated {N}-of-1 trials. Published online 2026.}

\bibitem[\citeproctext]{ref-sennSampleSizeConsiderations2019}
\CSLLeftMargin{6. }%
\CSLRightInline{Senn S. Sample size considerations for n-of-1 trials. \emph{Statistical Methods in Medical Research}. 2019;28(2):372-383. doi:\href{https://doi.org/10.1177/0962280217726801}{10.1177/0962280217726801}}

\bibitem[\citeproctext]{ref-zucker2010}
\CSLLeftMargin{7. }%
\CSLRightInline{Zucker DR, Ruthazer R, Schmid CH. Individual ({N-of-1}) trials can be combined to give population comparative treatment effect estimates: Methodologic considerations. \emph{Journal of Clinical Epidemiology}. 2010;63(12):1312-1323. doi:\href{https://doi.org/10.1016/j.jclinepi.2010.04.020}{10.1016/j.jclinepi.2010.04.020}}

\bibitem[\citeproctext]{ref-pmsimstats-paper10}
\CSLLeftMargin{8. }%
\CSLRightInline{{pmsimstats team}. Type {I} calibration of fixed-effect interaction tests in linear mixed models: A staged diagnosis of working-covariance misspecification and a cluster-robust remedy. Published online 2026.}

\bibitem[\citeproctext]{ref-pmsimstats-paper06}
\CSLLeftMargin{9. }%
\CSLRightInline{{pmsimstats team}. Three-component decomposition of treatment response in aggregated {N}-of-1 trials: Pharmacological, expectation-driven, and natural-history components. Published online 2026.}

\end{CSLReferences}

\end{document}
